Let $A$ and $A'$ be the original and edited lattices. \ui{Preserve Cartesian positions} determines how atoms adapt to the new cell. Equation~\eqref{eq:row-lattice-coordinates} gives
\begin{equation}
 \begin{aligned}
 \text{Keep fractional:}\quad&\mathbf f'=\mathbf f,
       &\mathbf r'&=\mathbf f A',\\
 \text{Keep Cartesian:}\quad&\mathbf r'=\mathbf r,
       &\mathbf f'&=\mathbf r(A')^{-1}.
 \end{aligned}
 \label{eq:lattice-preservation-options}
\end{equation}
For example, increase a cubic cell side from 4 to 8\,\AA. An atom initially at fractional position $(0.5,0,0)$ has Cartesian position $(2,0,0)\,\text{\AA}$. Without preserving Cartesian coordinates, it remains halfway along the cell and moves to $(4,0,0)\,\text{\AA}$. With the option enabled, it stays at $(2,0,0)\,\text{\AA}$ and its fractional position becomes $(0.25,0,0)$. Both choices can be useful; decide which geometric relationship to retain before editing the table.

\begin{table}[!htb]
\centering\small
\begin{tabularx}{\linewidth}{p{30mm}XX}
\toprule
Operation & Preserves & Main checks\\
\midrule
Edit lattice, keep fractional & $\mathbf f$ and site count & Cell lengths, angles, and changed interatomic distances.\\
Edit lattice, keep Cartesian & $\mathbf r$ and site count & New fractional positions and periodic boundary contacts.\\
Build supercell & Original periodic arrangement & Volume and independent site count multiplied by $|\det S|$.\\
Rotate lattice & Internal geometry and site count & Lattice and atoms rotating together.\\
Swap axes & Cartesian geometry and site count & Basis order and matching fractional components.\\
Mirror lattice & Fractional coordinates and site count & Reversed vector and cell handedness.\\
Apply strain & Fractional coordinates and site count & Vector lengths and affected distances.\\
\bottomrule
\end{tabularx}
\caption{Choose a lattice operation according to the quantities to preserve. Internal geometry means interatomic distances and angles; after changing periodic boundaries, inspect neighboring images in the new cell as well.}
\label{tab:round10-lattice-invariants}
\end{table}
